% Generated draft: requires mathematical and source review.
\documentclass[11pt]{article}
\usepackage[margin=25mm]{geometry}
\usepackage{fontspec}
\setmainfont{Latin Modern Roman}
\usepackage{amsmath,amssymb,mathtools,unicode-math}
\setmathfont{Latin Modern Math}
\usepackage{enumitem,needspace,xurl,hyperref}
\setlength{\parindent}{0pt}
\setlength{\parskip}{5pt}
\setlength{\emergencystretch}{3em}
\newcommand{\N}{\mathbb{N}}
\newcommand{\Z}{\mathbb{Z}}
\newcommand{\Q}{\mathbb{Q}}
\newcommand{\R}{\mathbb{R}}
\newcommand{\C}{\mathbb{C}}
\newcommand{\F}{\mathbb{F}}
\newcommand{\A}{\mathbb{A}}
\newcommand{\PP}{\mathbb{P}}
\newcommand{\E}{\mathbb{E}}
\newcommand{\eps}{\varepsilon}
\newcommand{\dd}{\,\mathrm{d}}
\begin{document}
\section*{UnsolvedMath ID: 701124}
\textbf{Source problem code:} catalogue entry
\subsection{Definitions and mathematical statement}
Work in the first-order language of groups \(\mathcal L_{\mathrm{grp}}=\{\cdot,{}^{-1},1\}\). A group quasi-identity is a sentence of the form
\[
\forall x_1\cdots\forall x_n\left(\bigwedge_{i=1}^{m}u_i(\mathbf x)=v_i(\mathbf x)\;\Longrightarrow\;u(\mathbf x)=v(\mathbf x)\right),
\]
where all displayed terms are group terms and \(m\geq 0\), with the empty conjunction interpreted as true. For a finite group \(F\), let \(\mathcal Q(F)\) be the quasivariety generated by \(F\), namely the class of all groups satisfying every group quasi-identity valid in \(F\).

For a group \(G\) and \(a\in G\), define
\[
a^G=\{g^{-1}ag:g\in G\},\qquad
\langle a^G\rangle=\langle g^{-1}ag:g\in G\rangle.
\]
Thus \(\langle a^G\rangle\) is the smallest normal subgroup of \(G\) containing \(a\). For any class \(M\) of groups, put
\[
L(M)=\{G:\text{for every }a\in G,\ \langle a^G\rangle\in M\}.
\]
\textbf{Question.} For every finite group \(F\), is \(L(\mathcal Q(F))\) finitely axiomatizable?

Under the usual first-order interpretation of finite axiomatizability, this asks whether, for every finite group \(F\), there exists a finite set \(\Sigma_F\) of \(\mathcal L_{\mathrm{grp}}\)-sentences such that, for every group \(G\),
\[
G\models\Sigma_F
\quad\Longleftrightarrow\quad
\forall a\in G\;\bigl(\langle a^G\rangle\in\mathcal Q(F)\bigr).
\]
The finite set may depend on \(F\), and \(G\) is not assumed finite. The supplied source does not specify whether the axioms must instead be quasi-identities. Under that more restrictive reading, the same question requires \(\Sigma_F\) to consist of quasi-identities; this qualification remains unresolved.
\subsection{Short English statement}
For each finite group \(F\), consider all groups in which the normal closure of every element satisfies every quasi-identity valid in \(F\). Must this class admit a finite axiomatization? The source leaves the permitted form of the axioms unspecified.
\subsection{Sources}
\begin{itemize}
\item Local UnsolvedMath registry, permanent ID \(701124\), exact source problem code \texttt{catalogue entry}, title \emph{Kourovka Notebook Problem 15.36}. The supplied record labels the problem \texttt{open}; this label has not been independently verified.
\item A. I. Budkin, proposer of Problem 15.36 in \emph{The Kourovka Notebook (Unsolved Problems in Group Theory)}, 15th issue (2002), as attributed in the supplied record.
\item The supplied background also identifies the 21st issue edition (2026), \texttt{arXiv:1401.0300v48}. The registry describes the Issues 1--21 collection as edited by E. I. Khukhro and V. D. Mazurov. These bibliographic details are supplied attribution and have not been independently checked.
\end{itemize}
\end{document}
